\begin{afterword}
\summaryhead

The post builds on the previous day's example, Japan's blood-type theory of personality. It invents a word, ``wiggin'', defined as a person with green eyes and black hair. In a short story, Danny sees such a man, calls him a wiggin, and expects him to commit crimes and reach for the ketchup. When his sister asks for evidence, he points to the dictionary.

The author explains that words carry us from visible features to inferred ones, and that dictionaries list the visible features used to recognize a thing, not the many properties people infer from it. So a connotation can attach to a word without appearing in its definition. Whenever someone pulls out a dictionary, the post says, they are generally trying to sneak in such a connotation, since the real question would take work to answer. The test is to restate the argument with only the definition in it (he is a wiggin, ``therefore, he's got black hair''). It sounds empty, so only the smuggled connotation makes such arguments feel like a point scored.

\responsehead

The post's central argument holds for the case it describes. If a definition lists only visible features, an argument from the definition can return only those features, and anything more has come from somewhere else. Danny shows this plainly. The test of restating his argument with only the definition in it is a check a reader can apply to any argument ``by definition''. The post's guess about how connotations arise also fits the real case it starts from: the Japanese belief about blood types is traced to the popular books of a journalist in the 1970s.

Some steps around that argument are weaker than the conclusion needs.

The question ``why not just call those people `green-eyed black-haired people'?'' is offered as if the existence of the word showed that it carries more than its definition. It does not. A short word for a phrase used often is useful in itself, as the author had written six days earlier. The next question, why not ask directly whether the man reaches for the ketchup, makes the point without this weakness.

The explanation ``people are lazy'' is a general claim about why people argue as they do, and the post gives no evidence for it.

The test also covers less than the last paragraphs claim. A valid argument from a definition ``would either say nothing new, or beg the question'': that is true of one-step arguments like Danny's. A longer deduction from definitions can tell a hearer something new, as the next post grants for mathematics. Where a rule attaches consequences to a definition, as in law, an argument from the definition can settle something real with no connotation involved.

So the thesis, that people who pull out a dictionary are ``generally'' sneaking in a connotation, is hedged and made plausible by the test. In this post its support is that test and one invented story.

In short: a sound test for arguments from a dictionary, with a rhetorical question that brevity answers, an unsupported claim that people are lazy, and a conclusion stated for all arguments from definitions that the test supports for one-step arguments like Danny's.
\end{afterword}
